\documentclass{article}%
\usepackage[T1]{fontenc}%
\usepackage[utf8]{inputenc}%
\usepackage{lmodern}%
\usepackage{textcomp}%
\usepackage{lastpage}%
\usepackage{amsmath}%
\usepackage{graphicx}%
%
\usepackage{graphicx}%
\title{Aérodynamique des bâtiments — Rapport complet}%
\author{FoamPilot}%
\date{\today}%
%
\begin{document}%
\normalsize%
\maketitle%
\tableofcontents%
\begin{abstract}%
Ce rapport présente la simulation aérodynamique urbaine autour d'un quartier de bâtiments avec simpleFoam et le modèle k{-}epsilon.%
\end{abstract}%
\section{Profil de couche limite urbaine}%
\label{sec:Profildecouchelimiteurbaine}%

%
\[%
u(y) = u_* \frac{\ln(y / y_0)}{\kappa}%
\]%
\section{Ratio d'aspect du canyon}%
\label{sec:Ratiodaspectducanyon}%

%
\[%
AR = \frac{H_{building}}{W_{street}} = 1.0%
\]%
\section{Paramètres de simulation}%
\label{sec:Paramtresdesimulation}%

%


\begin{figure}[h!]%
\begin{tabular}{|c|c|c|}%
\hline%
Paramètre&Valeur&Unité\\%
\hline%
Paramètre&Valeur&Unité\\%
\hline%
Vitesse entrée&10&m/s\\%
\hline%
Intensité turbulente&15&\%\\%
\hline%
Modèle turbulence&k{-}epsilon&\\%
\hline%
Temps final&100&s\\%
\hline%
Taille maille&40×20×10&cells\\%
\hline%
\end{tabular}%
\caption{Simulation parameters}%
\end{figure}

%
\section{Statistiques du maillage}%
\label{sec:Statistiquesdumaillage}%

%


\begin{figure}[h!]%
\begin{tabular}{|c|c|}%
\hline%
Statistic&Value\\%
\hline%
Statistique&Valeur\\%
\hline%
num\_points&11459\\%
\hline%
num\_cells&12786\\%
\hline%
bounds&{[}0.0, 200.0, 0.0, 100.00399780273438, 0.0, 50.0{]}\\%
\hline%
volume&1000000.0684243309\\%
\hline%
area&None\\%
\hline%
\end{tabular}%
\caption{Mesh quality statistics}%
\end{figure}

%
\section{Statistiques du champ de vitesse (U)}%
\label{sec:Statistiquesduchampdevitesse(U)}%

%


\begin{figure}[h!]%
\begin{tabular}{|c|c|}%
\hline%
Statistic&Value\\%
\hline%
Statistique&Valeur\\%
\hline%
mean&2.4451801776885986\\%
\hline%
std&4.262539863586426\\%
\hline%
min&{-}0.4652418792247772\\%
\hline%
max&12.615300178527832\\%
\hline%
volume\_weighted\_mean&{[}8.617223661765742, {-}5.279255594863886e{-}05, 0.000153282820345665{]}\\%
\hline%
\end{tabular}%
\caption{Velocity field (U) statistics}%
\end{figure}

%
\section{Statistiques du champ de pression (p)}%
\label{sec:Statistiquesduchampdepression(p)}%

%


\begin{figure}[h!]%
\begin{tabular}{|c|c|}%
\hline%
Statistic&Value\\%
\hline%
Statistique&Valeur\\%
\hline%
mean&0.5063239336013794\\%
\hline%
std&1.0810129642486572\\%
\hline%
min&{-}4.667999744415283\\%
\hline%
max&8.727890014648438\\%
\hline%
volume\_weighted\_mean&0.7527323818939506\\%
\hline%
\end{tabular}%
\caption{Pressure field (p) statistics}%
\end{figure}

%
\section{Visualisations}%
\label{sec:Visualisations}%

%


\begin{figure}[h!]%
\centering%
\includegraphics[width=0.7\textwidth]{/home/steven/foampilot/foampilot/tutorials/06_buildingAero/slice_plot.png}%
\caption{Slice Plot.Png}%
\end{figure}

%


\begin{figure}[h!]%
\centering%
\includegraphics[width=0.7\textwidth]{/home/steven/foampilot/foampilot/tutorials/06_buildingAero/contour_plot.png}%
\caption{Contour Plot.Png}%
\end{figure}

%


\begin{figure}[h!]%
\centering%
\includegraphics[width=0.7\textwidth]{/home/steven/foampilot/foampilot/tutorials/06_buildingAero/vector_plot.png}%
\caption{Vector Plot.Png}%
\end{figure}

%


\begin{figure}[h!]%
\centering%
\includegraphics[width=0.7\textwidth]{/home/steven/foampilot/foampilot/tutorials/06_buildingAero/mesh_style_plot.png}%
\caption{Mesh Style Plot.Png}%
\end{figure}

%
\appendix%
\section{Données cellule}%
\label{sec:Donnescellule}%
Les données cellule ont été exportées vers cell\_data.csv pour analyse externe.

%
\end{document}